\section{Causal spectral synthesis for the actual weighted Paley--Wiener levels}
\label{sec:global-weighted-pw}

The polynomial weight cannot be removed by the equivalent-norm
argument of MA-GP-003.  We instead identify the actual theta range
as the closure of the range of a bounded causal convolution.
An explicit divided-difference identity gives the quotient
resolvent.  Its principal parts, polynomial bounds, and spectral
completeness are proved below.  This supplies the half-line
Hardy--Sobolev closure without invoking a weighted analogue
of Beurling's theorem.

Fix $0<\alpha<1/2$, $s\ge0$, $Y>0$, and set
\[
 c=-\log Y,\qquad d=\tfrac12-\alpha,\qquad
 u=\tfrac12+\alpha=1-d,\qquad
 \mathcal H_s=L^2((0,\infty),(1+y^2)^{s/2}\,dy).
 \tag{MA-GW.1}
\]
The actual global space is
$H=L^2(\mathbb R,e^{2\alpha|x|}(1+x^2)^{s/2}dx)$,
and $H_c$ is its supported half-line subspace.
The map $U_cf(y)=e^{\alpha(c+y)}f(c+y)$ carries its actual
norm to
\[
 \|U_cf\|_{W_c}^2,\qquad
 W_c(y)=e^{2\alpha|c+y|-2\alpha(c+y)}
                              (1+(c+y)^2)^{s/2}.
 \tag{MA-GW.2}
\]
For this fixed $c$ the $W_c$ norm and the $\mathcal H_s$
norm are equivalent.  The actual closed subspace $J_Y$
is always the theta-image closure from MA-GP-001--002;
it is not replaced by an unproved weighted inner-function
range.

\paragraph{MA-GW-001 (an actual bounded causal range).}
For a sufficiently large integer $N$ define
\[
               \Phi(z)=\frac{(z-u)\zeta(z+d)}{(z+2)^N}.
 \tag{MA-GW.3}
\]
The singularity at $z=u$ is removable and its value is nonzero.
There is a continuous causal kernel $\phi$, supported in
$[0,\infty)$, with $|\phi(t)|\le C e^{-\eta t}$ for some
$\eta>0$, whose Laplace transform is $\Phi$.
Consequently
\[
 T_\phi f=\int_0^\infty\phi(t)S_tf\,dt,\qquad
 (S_tf)(y)=\mathbf1_{y\ge t}f(y-t),
 \tag{MA-GW.4}
\]
is bounded on $\mathcal H_s$, and the actual range satisfies
\[
                  U_cJ_Y=\overline{\operatorname{ran}T_\phi}
                                           ^{\,\mathcal H_s}.
 \tag{MA-GW.5}
\]
The same equality holds as a set in the equivalent $W_c$ norm.

\begin{proof}
\emph{The kernel and the weighted semigroup.}
Take $0<\eta<d$, for example $\eta=d/2$.
The numerator of (MA-GW.3) is holomorphic on the whole plane,
including $z=u$, because the residue of zeta at $1$ is one.
Its denominator has no zero on $\operatorname{Re}z\ge-\eta$.
The unconditional fixed-strip polynomial bound for zeta
from NT-Z10, and the Euler bound on a right half-plane,
allow $N$ to be chosen so that
\[
 |\Phi(\sigma+it)|\le C(1+|t|)^{-2}
                       \quad(-\eta\le\sigma\le2),
 \tag{MA-GW.6}
\]
with the removable value included on bounded sets.
For $\sigma\ge2$, its Euler product also gives the uniform
bound $|\Phi(\sigma+it)|\le C|\sigma+it|^{1-N}$.
Increase $N$ if necessary so $N>3$.

Define the inverse Laplace integral on a vertical line with
positive real part.  It is absolutely convergent by
(MA-GW.6).  Cauchy's theorem, first on finite rectangles
and then by the uniform vertical decay, allows its line
to move to $\operatorname{Re}z=-\eta$.
This proves the continuous-kernel bound
$|\phi(t)|\le C e^{-\eta t}$ for $t\ge0$.
For $t<0$, move the line to $\operatorname{Re}z=M$ and
then let $M\to\infty$.
The absolute integral there is at most
$C e^{Mt}M^{2-N}$, so it tends to zero.
Thus the kernel is causal.  Fourier inversion on a
positive vertical line proves that its Laplace transform
is exactly $\Phi$; the identity then holds throughout
$\operatorname{Re}z>-\eta$ by analytic continuation.

The shifts are strongly continuous on $\mathcal H_s$ and obey
\[
 \|S_t\|\le2^{s/4}(1+t^2)^{s/4}\qquad(t\ge0).
 \tag{MA-GW.7}
\]
The inequality follows from
$1+(y+t)^2\le2(1+y^2)(1+t^2)$; continuity follows first
on compact smooth functions and then by density.
The exponential kernel bound therefore makes (MA-GW.4)
a bounded vectorwise integral.  Its Laplace transform
is multiplication by $\Phi$, by Fubini on a right
half-plane and then uniqueness.

\emph{The exact source primitive.}
Let $v$ be an even smooth frequency test supported in
$[-Y,Y]$ with $v(0)=0$ and $\int v=0$.
Write $f=\mathcal F_e v$ for the even Fourier inverse.
The actual Poisson theta formula gives, with
$z$ initially far enough to the right,
\[
 (\mathcal L U_c\Theta f)(z)
       =2e^{zc}\zeta(z+d)
                        \int_0^Yv(t)t^{z+d-1}\,dt .
 \tag{MA-GW.8}
\]
Both sides continue to the open right half-plane.
Make the change $t=Y e^{-q}$ and put
\[
 a(q)=e^{-dq}v(Y e^{-q}),\qquad
 b(q)=e^{uq}\int_0^q e^{-ur}a(r)\,dr.
 \tag{MA-GW.9}
\]
The moment condition gives
\[
       \int_0^\infty e^{-uq}a(q)\,dq
                         =Y^{-1}\int_0^Yv(t)\,dt=0.
 \tag{MA-GW.10}
\]
Thus $a=b'-ub$ and also
$b(q)=-e^{uq}\int_q^\infty e^{-ur}a(r)\,dr$.
At zero, $a$ and $b$ are flat because $v$ is smooth and
flat at the band endpoint $Y$.
At infinity, evenness and $v(0)=0$ give
$v(t)=O(t^2)$, including the corresponding logarithmic
derivative bounds.  Hence $a,b$ and every derivative
are $O(e^{-(d+2)q})$.
In particular $b$ and all its derivatives belong to
$\mathcal H_s$, with all zero boundary traces at zero.

If $B(z)=\mathcal Lb(z)$, (MA-GW.8)--(MA-GW.10) become
the exact identity
\[
             (\mathcal L U_c\Theta f)(z)
                 =2Y^d(z-u)\zeta(z+d)B(z)
                 =2Y^d\Phi(z)(z+2)^N B(z).
 \tag{MA-GW.11}
\]
The factor $2Y^d$ is fixed and nonzero.
Since $(\partial+2)^N b\in\mathcal H_s$ and all its
integration-by-parts boundary terms vanish, (MA-GW.11)
shows that every such actual theta core vector belongs
to $\operatorname{ran}T_\phi$.
The band-supported Schwartz tests used to define the
source range have exactly this smoothness and flatness;
equivalently one may use the inward-dilation approximation
already justified in MA-GP-005.

Conversely, every $b\in C_c^\infty((0,\infty))$ gives a
source test by setting $a=b'-ub$ and
$v(Y e^{-q})=e^{dq}a(q)$, then extending evenly.
The result is smooth, supported away from both zero and
the band endpoints, and has zero integral by (MA-GW.10).
Thus the vector with Laplace transform
$(z-u)\zeta(z+d)B(z)$ is in the actual theta core.

It remains to approximate all of $\operatorname{ran}T_\phi$.
Take $h\in C_c^\infty((0,\infty))$ and solve
$(\partial+2)^N b=h$ with zero initial traces by the
causal kernel $e^{-2q}q^{N-1}/(N-1)!$.
The resulting $b$ is flat at zero and has exponentially
decreasing derivatives at infinity.
Smooth compact truncations $b_M$ therefore satisfy
$(\partial+2)^N b_M\to h$ in $\mathcal H_s$.
Boundedness of $T_\phi$ and the preceding source construction
show that $T_\phi h$ is in $U_cJ_Y$.
Such $h$ are dense in $\mathcal H_s$.
This proves both inclusions in (MA-GW.5).
Equivalence of the two fixed-$c$ norms proves the final assertion.
\end{proof}

\paragraph{MA-GW-002 (weighted half-line spectral synthesis and trace).}
Let
\[
 b_s=\#\{j\in\mathbb Z_{\ge0}:2j<s-1\},\qquad
 r^+_\rho=
 \begin{cases}
 m_\rho,&\operatorname{Re}\rho<u,\\
 \min(m_\rho,b_s),&\operatorname{Re}\rho=u,\\
 0,&\operatorname{Re}\rho>u.
 \end{cases}
 \tag{MA-GW.12}
\]
The exact actual closure in $H_c$ is the common kernel
of the continuous Mellin jets of orders
$0\le j<r^+_\rho$ at $\rho-1/2$.
The quotient $\mathcal Q_c=H_c/J_Y$ has dense span of
its finite spectral blocks, with precisely these lengths.
For every additive $g\in C_c^\infty((0,\infty))$ its
actual quotient convolution is trace class and
\[
       \operatorname{Tr}_{\mathcal Q_c}C_g^{\mathcal Q_c}
           =\sum_\rho r^+_\rho\,\mathcal Mg(\rho-1/2).
 \tag{MA-GW.13}
\]
The sum converges absolutely and is independent of $Y$.
A sufficiently high power of the actual quotient
semigroup resolvent is trace class.
No zero separation or positive distance from the upper
edge is assumed.

\begin{proof}
\emph{The quotient and its meromorphic resolvent are actual operators.}
Work first on $\mathcal H_s$ and let
$\mathcal J=\overline{\operatorname{ran}T_\phi}$.
The shifts commute with $T_\phi$, so they preserve
$\mathcal J$ and induce a strongly continuous semigroup
on $\mathcal Q=\mathcal H_s/\mathcal J$ with the same
polynomial bound (MA-GW.7).
Let $A$ be its generator.
For $\operatorname{Re}\lambda>0$ the vectorwise
Laplace integral gives its bounded resolvent
\[
          R(\lambda)=(\lambda-A)^{-1}
                     =\int_0^\infty e^{-\lambda t}S_t^{\mathcal Q}\,dt .
 \tag{MA-GW.14}
\]
The direct difference-quotient and generator-domain
integration-by-parts verification is the one given in
MA-GP-009; (MA-GW.7) supplies every integrable majorant.

For $\operatorname{Re}\lambda<\eta$ define the scalar kernel
and actual bounded operator
\[
 \psi_\lambda(r)=\int_r^\infty\phi(t)e^{\lambda(t-r)}\,dt,
 \qquad
 T_\lambda=\int_0^\infty\psi_\lambda(r)S_r^{\mathcal Q}\,dr .
 \tag{MA-GW.15}
\]
The map $\lambda\mapsto\psi_\lambda$ is holomorphic in
the weighted $L^1$ kernel norm used with (MA-GW.7).
Differentiating its scalar exponential and using
$|\phi(t)|\le Ce^{-\eta t}$ proves this on each smaller
half-plane.  Hence $T_\lambda$ is operator-norm
holomorphic there.  Moreover
\[
 \sup_{\operatorname{Re}\lambda\le\epsilon}\|T_\lambda\|<\infty
                      \qquad(\epsilon<\eta).
 \tag{MA-GW.16}
\]
For example, bound its double integral by
$C_s\int_0^\infty|\phi(t)|e^{\max(\epsilon,0)t}
\int_0^t(1+r)^{s/2}\,dr\,dt$.
No bound here depends on the imaginary part or on how
far $\operatorname{Re}\lambda$ is negative.

The derivative $\psi_\lambda'=-\phi-\lambda\psi_\lambda$
and the value $\psi_\lambda(0)=\Phi(-\lambda)$ give
the bounded divided-difference identity
\[
 (\lambda-A)T_\lambda=\Phi(-\lambda)\,1,\qquad
 T_\lambda(\lambda-A)=\Phi(-\lambda)\,1
                                      \quad\hbox{on }\mathcal D(A).
 \tag{MA-GW.17}
\]
For clarity, before taking the quotient the right side is
$\Phi(-\lambda)1-T_\phi$.
Integrating the semigroup derivative against
$\psi_\lambda$ gives this formula, with boundary value
$\psi_\lambda(0)$; a shift difference quotient proves
that the resulting operator maps into the generator
domain.  On the quotient $T_\phi$ induces the zero
operator by (MA-GW.5), proving (MA-GW.17).
All these convolution operators preserve the range
closure because they commute with $T_\phi$.

Whenever $\Phi(-\lambda)\ne0$, the two identities show that
\[
                         R(\lambda)=T_\lambda/\Phi(-\lambda).
 \tag{MA-GW.18}
\]
This is a two-sided bounded inverse of the actual closed
generator, not just a formal meromorphic function.
It agrees with (MA-GW.14) on their common domain.
Together they give a meromorphic continuation of the
resolvent to the plane.  In the right half-plane any
apparent singularity of (MA-GW.18) is removable by
(MA-GW.14).  The remaining possible poles are
\[
                \lambda_\rho=\rho-u,\qquad
                \operatorname{Re}\rho\le u,
 \tag{MA-GW.19}
\]
of order at most $m_\rho$.
Indeed the zeros of $\Phi(-\lambda)$ there correspond
to $\zeta(d-\lambda)=0$, and the functional equation
pairs this multiset with (MA-GW.19).
The factor at $d-\lambda=1$ is removable and nonzero.
Trivial zeta zeros would give $\operatorname{Re}\lambda>0$
and contribute no singularity to the actual resolvent.
Thus all possible poles have real parts in $[-u,0]$.

\emph{Every principal part has finite rank, with the exact edge truncation.}
Let $\lambda_0$ be one possible pole, with zero multiplicity $m$.
The contour projection
$E_{\lambda_0}=(2\pi i)^{-1}\int_\gamma R(\lambda)\,d\lambda$
is a bounded projection.  Here $\gamma$ is a small
positively oriented circle containing no other possible pole.
These familiar resolvent facts follow directly in the present
setting: integrate
$R(z)R(w)=(R(w)-R(z))/(z-w)$ over nested circles to obtain
$E_{\lambda_0}^2=E_{\lambda_0}$ and commutation with the
resolvent.  The identity $AR(z)=zR(z)-1$ allows integration
in the graph norm and gives
$A^jE_{\lambda_0}=(2\pi i)^{-1}\int_\gamma z^jR(z)\,dz$.
Since the possible pole order is at most $m$,
\[
                  (A-\lambda_0)^mE_{\lambda_0}=0.
 \tag{MA-GW.20}
\]
More explicitly, the coefficient of
$(z-\lambda_0)^{-j}$ in the principal part is
$(A-\lambda_0)^{j-1}E_{\lambda_0}$, by the same
graph-norm contour identity.
The same statements hold on the continuous linear dual.

Identify that dual with functions $v$ such that
\[
       \ell_v(f)=\int_0^\infty f(y)v(y)\,dy,\qquad
       \int_0^\infty |v(y)|^2(1+y^2)^{-s/2}\,dy<\infty.
 \tag{MA-GW.21}
\]
The dual of the quotient is the annihilator of $\mathcal J$.
The transpose semigroup acts by $v(y)\mapsto v(y+t)$.
Therefore a vector in the dual range of $E_{\lambda_0}$
satisfies $(\partial-\lambda_0)^m v=0$ on $(0,\infty)$
in the distributional sense.
Its complete solution space consists of
$e^{\lambda_0y}p(y)$ with $\deg p<m$:
multiply by $e^{-\lambda_0y}$ and successively integrate
the equation $\partial^m q=0$ on the connected interval.
This conclusion requires no diagonalizability assumption.

If $\operatorname{Re}\lambda_0<0$, all these $m$
polynomials meet (MA-GW.21).
If $\operatorname{Re}\lambda_0=0$, the polynomial of
highest degree $j$ meets it exactly when $2j<s-1$.
This gives the upper bounds $m$ and $\min(m,b_s)$
for the projection rank.
They are also lower bounds.
The functionals
\[
                         \ell_j(f)=\int_0^\infty
                                      y^j e^{\lambda_0y}f(y)\,dy
 \tag{MA-GW.22}
\]
for precisely these allowed $j<m$ are continuous and
annihilate $\operatorname{ran}T_\phi$.
To see the latter, substitute the causal convolution
and expand $(y+t)^j$: every coefficient is a derivative
of $\Phi(-\lambda)$ at $\lambda=\lambda_0$ of order
at most $j<m$, and is zero.
Absolute convergence follows from the exponential
kernel bound and (MA-GW.21).
Their translates form the usual finite triangular
Jordan chain.  The resolvent integral (MA-GW.14),
applied to these functionals, is the finite geometric
resolvent of that chain; its small contour integral
acts as the identity.  Thus they belong to the dual
range of $E_{\lambda_0}$.
They are linearly independent by their polynomial
densities.  This proves equality of the rank bounds,
and hence finite rank of every principal part.
It also identifies its full dual range as (MA-GW.22).

Under $U_c^{-1}$ the functionals (MA-GW.22) become,
up to a nonzero scalar,
\[
              f\longmapsto\int_c^\infty
                     (x-c)^j e^{(\lambda_0+\alpha)x}f(x)\,dx .
 \tag{MA-GW.23}
\]
Since $\lambda_0+\alpha=\rho-1/2$, their spans are
the ordinary Mellin jets at that point, by the
invertible triangular change from $(x-c)^j$ to $x^j$.
This proves the local multiplicity assertion (MA-GW.12).

\emph{Uniform polynomial contour bounds.}
Choose $0<\epsilon<\eta$ and $C_0>u+2$.
On $\operatorname{Re}\lambda=\epsilon$,
(MA-GW.14) and (MA-GW.7) bound $R(\lambda)$ uniformly
in its imaginary part.
On the entire region $\operatorname{Re}\lambda\le-C_0$,
the Euler lower bound for $\zeta(d-\lambda)$ and
(MA-GW.16)--(MA-GW.18) give the uniform estimate
\[
                 \|R(\lambda)\|\le C(1+|\lambda|)^{N-1}.
 \tag{MA-GW.24}
\]
Here $|-\lambda-u|\ge c(1+|\lambda|)$ after increasing
$C_0$, and the Euler lower bound is uniform for
$\operatorname{Re}(d-\lambda)\ge d+C_0>1$.
In particular the constant in (MA-GW.24) does not grow
as the left line is moved farther left.

Use the good heights of MT-M18, (M18.7)--(M18.8), reflected
if necessary, to choose increasing $t_n=n+O(1)$ with
$1/2<t_{n+1}-t_n<3/2$.
That full minimum-modulus proof supplies a polynomial
lower bound for the finite zeta function on
$[-C_0,\epsilon]+it_n$ after the change $z=d-\lambda$.
It is independent of zero separation inside each block.
Combining it with (MA-GW.16) and the explicit polynomial
factors in $\Phi$ gives constants $C,B$ such that
\[
 \sup_{\lambda\in\Gamma_n}\|R(\lambda)\|
      \le C(2+|n|)^B,\qquad
 \Gamma_n=\partial\bigl([-C_0,\epsilon]
                              +i[t_n,t_{n+1}]\bigr).
 \tag{MA-GW.25}
\]
The bound on the right side uses the actual resolvent
integral, even if an apparent zero of the denominator
there has been canceled.
The horizontal sides avoid all zeros by construction.
All low-height contours can be adjusted individually.
Their lengths are uniformly bounded.
Their projections $E_n$ have finite rank
\[
                           \operatorname{rank}E_n
                                     =O(\log(2+|n|)),
 \tag{MA-GW.26}
\]
by the finite-rank proof and NT-Z10.
These estimates use one fixed central strip, together
with the globally uniform left-region bound
(MA-GW.24); no unstated uniform minimum modulus on
expanding strips is needed.

\emph{Causal smoothing removes every residual subspace.}
For $k\in C_c^\infty((0,\infty))$ put
$H_k(\lambda)=\int_0^\infty k(t)e^{\lambda t}\,dt$.
It decreases faster than every power of
$|\operatorname{Im}\lambda|$ on the fixed central strip.
Therefore the finite-rank operators
\[
       V_n(k)=\frac1{2\pi i}\int_{\Gamma_n}
                                  H_k(\lambda)R(\lambda)\,d\lambda
 \tag{MA-GW.27}
\]
have an absolutely convergent sum in trace norm:
\[
 \sum_n\|V_n(k)\|_1
 \le C_M\sum_n\log(2+|n|)(2+|n|)^{B-M}<\infty
                                             \quad(M>B+2).
 \tag{MA-GW.28}
\]
Their ranges are contained in $\operatorname{ran}E_n$,
by the finite-dimensional resolvent calculus already
proved from the contour identities.

To identify the sum as an actual operator, first insert
(MA-GW.14) on the line $\operatorname{Re}\lambda=\epsilon$.
Fourier inversion and Fubini, justified by the rapid
decay of $H_k$ and the exponentially damped polynomial
semigroup bound, give
\[
 C_k^{\mathcal Q}
     =\frac1{2\pi i}\int_{\epsilon-i\infty}^{\epsilon+i\infty}
                                    H_k(\lambda)R(\lambda)\,d\lambda .
 \tag{MA-GW.29}
\]
Summing finite rectangles and letting their good heights
tend to both infinities moves this integral to the line
$-C_0$ and picks up exactly (MA-GW.27).
The horizontal integrals tend to zero by (MA-GW.25).
One may then move the left line to $-M$ for any $M>C_0$:
there are no poles to its left, and (MA-GW.24) controls
the added horizontal segments.

Finally let $M\to\infty$.
If $\operatorname{supp}k\subset[a,b]$ with $a>0$,
integration by parts gives, for every integer $\ell$,
\[
 |H_k(-M+it)|
    \le C_{\ell,k}e^{-Ma}(1+M)^\ell(1+|t|)^{-\ell}.
 \tag{MA-GW.30}
\]
Combining (MA-GW.24) with (MA-GW.30), for $\ell>N+1$,
makes the norm of the remaining vertical integral tend
to zero, bounded by an exponentially decreasing factor
times a fixed power of $1+M$.
Consequently
\[
                    C_k^{\mathcal Q}=\sum_n V_n(k)
                              \quad\hbox{in trace norm}.
 \tag{MA-GW.31}
\]
This establishes an actual trace-class operator identity,
not just equality of its possible residues.

Choose a smooth positive-time approximate identity
$k_\varepsilon$ of integral one with support in
$[\varepsilon,2\varepsilon]$.
Strong continuity of the semigroup gives
$C_{k_\varepsilon}^{\mathcal Q}\to1$ strongly.
Equation (MA-GW.31) then proves density of the span
of the finite spectral blocks.
It also proves the exact closure assertion:
if every functional (MA-GW.22) vanishes on a quotient
vector, every $E_n$ kills that vector, hence
every $V_n(k)$ and every $C_k^{\mathcal Q}$ kill it,
and the approximate identity makes the vector zero.
There is no remaining subspace invisible to the
local principal parts.

On each finite block, the trace in (MA-GW.27) is its
length times $H_k(\lambda_\rho)$; the higher Jordan
terms have trace zero.
Trace-norm convergence permits summation of those
finite traces.
For the physical test $g$, take $k(t)=e^{\alpha t}g(t)$.
Then $H_k(\lambda_\rho)=\mathcal Mg(\rho-1/2)$.
Bounded similarity back to the actual $H_c/J_Y$
proves (MA-GW.13), including its absolute convergence.

\emph{A trace-class resolvent power.}
Fix $\nu>\epsilon$ and choose an integer
$q>\max\{B+2,N+1\}$.
Repeat the same contours with
$H(\lambda)=(\nu-\lambda)^{-q}$.
The right-line integral is
$(\nu-A)^{-q}$, by its vectorwise Laplace representation
with kernel $t^{q-1}e^{-\nu t}/(q-1)!$ and Fourier inversion.
The trace-norm block series converges by (MA-GW.25)--(MA-GW.26).
On the line $-M$, (MA-GW.24) bounds the remaining
integral by $C M^{N-q}$, which tends to zero.
The central horizontal integrals vanish by $q>B+2$,
and the farther-left ones by $q>N+1$.
Thus
\[
                              (\nu-A)^{-q}\in\mathcal S_1(\mathcal Q).
 \tag{MA-GW.32}
\]
Bounded similarity gives the same assertion for the
actual normalized quotient semigroup.
\end{proof}

\paragraph{MA-GW-003 (the actual weighted lift and the two orders of limits).}
Let $P_c^Z$ be the actual orthogonal projection onto
$Z_c=H_c\ominus J_Y$, extended by zero on $H_c^\perp$.
For every $g\in C_c^\infty((0,\infty))$,
\[
 C_gP_c^Z\in\mathcal S_1(H),\qquad
 \operatorname{Tr}_H(C_gP_c^Z)
        =\sum_\rho r^+_\rho\,\mathcal Mg(\rho-1/2).
 \tag{MA-GW.33}
\]
This is the actual convolution-after-projection operator.
For the fixed smooth exhaustion $\chi_R$ of MA-GP-009,
\[
 \begin{aligned}
 \lim_{Y\to\infty}\lim_{R\to\infty}
    \operatorname{Tr}(M_{\chi_R}C_gQ_Y^{\rm glob})
      &=-\sum_\rho r^+_\rho\,\mathcal Mg(\rho-1/2),\\
 \lim_{R\to\infty}\lim_{Y\to\infty}
    \operatorname{Tr}(M_{\chi_R}C_gQ_Y^{\rm glob})
      &=-\sum_\rho r^{\rm strip}_\rho\,\mathcal Mg(\rho-1/2),
 \end{aligned}
 \tag{MA-GW.34}
\]
where the actual full-strip multiplicities are
\[
 r^{\rm strip}_\rho=
 \begin{cases}
 m_\rho,&d<\operatorname{Re}\rho<u,\\
 \min(m_\rho,b_s),&\operatorname{Re}\rho=d
                                \ \hbox{or}\ \operatorname{Re}\rho=u,\\
 0,&\hbox{otherwise}.
 \end{cases}
 \tag{MA-GW.35}
\]
The first limit minus the second is exactly
\[
 -\sum_{\operatorname{Re}\rho<d}m_\rho\,\mathcal Mg(\rho-1/2)
 -\sum_{\operatorname{Re}\rho=d}
          \bigl(m_\rho-\min(m_\rho,b_s)\bigr)
                                    \mathcal Mg(\rho-1/2).
 \tag{MA-GW.36}
\]
Thus all $s\ge0$, including the draft's $s\ge3$, are covered.
No assertion that an omitted zero or omitted edge jet actually
exists is part of this formula.

\begin{proof}
\emph{The weighted adjoint domain on the actual half-line.}
Use the actual unitary $U_c:H_c\to H_{W_c}$ from
(MA-GW.2), and write $W=W_c$.
The normalized unilateral semigroup has adjoint
\[
 (S_t^{*,W}v)(y)=\frac{W(y+t)}{W(y)}v(y+t),
            \qquad L_Wv=W^{-1}\partial_y(Wv).
 \tag{MA-GW.37}
\]
The orthogonal complement of the actual ideal is invariant
under these adjoints.
Multiplication by $W$ is unitary onto $L^2(W^{-1}dy)$,
where the adjoints become left translations.
Their generator has no boundary condition at zero.
Exactly as in the explicit domain proof of MA-GP-009,
\[
 v\in\mathcal D(L_W^q)
 \quad\Longleftrightarrow\quad
 \partial^j(Wv)\in L^2(W^{-1}dy)\quad(0\le j\le q).
 \tag{MA-GW.38}
\]
The proof uses the fundamental theorem of calculus for
the direct difference quotient and interior test functions
for its converse, followed by induction.
The polynomial shift bound and fixed-$c$ norm equivalence
give the required translation bounds.
In particular $Wv\in H^q_{\rm loc}([0,\infty))$.

The trace-class quotient resolvent power of (MA-GW.32),
bounded similarity to the actual quotient, and Hilbert
adjunction give
\[
 \mathcal R_q=(\nu-L_W)^{-q}P_{Z_W}
                         \in\mathcal S_1(H_W),
                 \qquad Z_W=U_cZ_c .
 \tag{MA-GW.39}
\]
The full resolvent restricts to the adjoint invariant
subspace by its vectorwise Laplace integral.
The identities
$L_W(\nu-L_W)^{-1}=\nu(\nu-L_W)^{-1}-1$
show that $\mathcal R_q$ maps boundedly into the
$q$th graph domain.

\emph{All endpoint and cusp terms are finite rank.}
Let $k(t)=e^{\alpha t}g(t)$, extended by zero outside
$(0,\infty)$, and take $v\in\mathcal D(L_W^q)$.
Put $b=Wv$ and $a=W^{-1}$.
The exact integral before integration by parts is
\[
 C_k^+(\nu-L_W)^qv(x)
       =\int_0^\infty k(x-y)a(y)
                              (\nu-\partial_y)^q b(y)\,dy .
 \tag{MA-GW.40}
\]
When $c<0$, split this integral at $L=-c$; otherwise
there is no interior split.
On either smooth part all functions
$W\partial^j(W^{-1})$ are bounded.
Indeed
\[
 \log W(y)=2\alpha|c+y|-2\alpha(c+y)
                              +\frac{s}{2}\log(1+(c+y)^2),
 \tag{MA-GW.41}
\]
whose positive-order derivatives are bounded on each
part and have finite one-sided values at the cusp.
The bulk after $q$ integrations by parts is
\[
 \mathcal B_{k,q}v(x)
  =\int W(y)(\nu+\partial_y)^q
                         \bigl(k(x-y)/W(y)\bigr)v(y)\,dy .
 \tag{MA-GW.42}
\]
It is a finite sum of bounded causal convolutions
$C_{k^{(i)}}^+$ followed by bounded input multipliers.
Hence it is bounded on $H_W$.

For explicit control of the boundary signs and orders,
the term of derivative order $j\ge1$ on each interval
has boundary expression
\[
 \binom qj\nu^{q-j}
       \sum_{\ell=0}^{j-1}(-1)^{j+\ell}
        \left[\partial_y^\ell(k(x-y)a(y))\,
                    b^{(j-1-\ell)}(y)\right]_{\rm left}^{\rm right}.
 \tag{MA-GW.43}
\]
Terms at infinity vanish by the compact support of $k$.
At zero take the lower-endpoint term, and at $L$ take
the difference of the two one-sided expressions.
Because $b\in H^q_{\rm loc}$, its derivatives through
order $q-1$ have common traces at the cusp.
Every remaining output profile is a finite linear
combination of $k^{(i)}(x)$ or $k^{(i)}(x-L)$.
It is smooth and compactly supported in the half-line.
Every input functional is a trace $b^{(j)}(0)$ or
$b^{(j)}(L)$ with $j<q$.
The elementary one-dimensional Sobolev trace inequality
bounds these functionals by the graph norm in (MA-GW.38).
For $q=2$ and $s=0$, (MA-GW.43) specializes exactly
to the endpoint and $-4\alpha$ cusp terms of
(MA-GP.102)--(MA-GP.104).

Apply (MA-GW.40)--(MA-GW.43) to $v=\mathcal R_q f$.
Its left side is $C_k^+P_{Z_W}f$.
Its right side is the product
$\mathcal B_{k,q}\mathcal R_q$ plus finitely many
rank-one operators, since each graph trace composed
with $\mathcal R_q$ is bounded.
This proves actual ambient trace class.
Conjugating by $U_c$ and extending by zero gives
$C_gP_c^Z\in\mathcal S_1(H)$.
Causality ensures that $C_g$ preserves $H_c$.
Trace cyclicity with the actual orthogonal projection
therefore identifies its trace with quotient convolution,
already computed in (MA-GW.13).
This proves (MA-GW.33).

\emph{Both actual iterated limits.}
Let $P_c=\mathbf1_{[c,\infty)}$ on $H$.
The orthogonal decomposition is still
$Q_Y^{\rm glob}=P_c-P_c^Z$.
For compact smooth $\chi$, $M_\chi C_gP_c$ is trace
class with trace zero.
The proof is the one given in MA-GP-009: restrict a
smooth compact convolution kernel by bounded position
multipliers, conjugate by the two bounded compact
coordinate factors of the weight, and approximate
the half-line indicator strongly by smooth multipliers.
The diagonal is $g(0)\chi(x)=0$.
All these steps apply to the present polynomial weight
because it is bounded above and below on each compact
coordinate projection.
Thus
\[
 \operatorname{Tr}(M_{\chi_R}C_gQ_Y^{\rm glob})
                    =-\operatorname{Tr}(M_{\chi_R}C_gP_c^Z).
 \tag{MA-GW.44}
\]
For fixed $Y$, strong convergence of the uniformly
bounded $M_{\chi_R}$ and (MA-GW.33) give convergence
in trace norm of the right-hand product to $C_gP_c^Z$.
Its trace is independent of $Y$, proving the first
line of (MA-GW.34).

For the opposite order, MA-GP-001 gives the actual
fixed-position trace-norm condensation for every
$s\ge0$.  The subsequent actual global cutoff limit
is MA-WP-003, with its exact two boundary-jet counts.
That is precisely (MA-GW.35), proving the second line.
All sums are absolutely convergent by NT-Z10 and
rapid vertical decay of the test transform.
Their difference is (MA-GW.36); the upper-edge terms
cancel, while only the unretained lower-edge jets
remain there.
\end{proof}

\paragraph{Scope, primary inputs, and trace interpretation.}
The argument uses the exact source theta and Poisson formulas,
the classical containing strip and unconditional fixed-strip
growth for zeta, and the complete minimum-modulus construction
of MT-M18.  Its required lower bound is (M18.7)--(M18.8)
for the finite zeta factor on one fixed strip, with the
disk-Blaschke/Harnack proof given there.
The analytic antecedent and local multiplicity count are
NT-Z10, with the Ramakrishnan--Valenza locators recorded
in that module.  All new causal range, resolvent, finite-rank,
contour, and actual lift arguments are supplied above.
Neither $\mathsf A$ nor $M$ is used, and no stronger
zero-location assertion than the classical containing
strip is an input.

The result supplies the precise polynomial-weight extension
that an ordinary $s=0$ Hardy identification cannot supply.
At the critical equalities $2j=s-1$, the norm in
(MA-GW.21) diverges logarithmically, so those jets are
excluded without an unsupported density assertion at the
Sobolev threshold.
No uniform trace-norm bound as $Y\to\infty$, no unspecified
interchange of joint limits, and no equality with a local
additive Fourier projection is inferred.

In the inherited Grothendieck trace framework, the actual
finite-length modules arise as the finite-rank principal
parts of the quotient resolvent, not as a presupposed
spectral basis.  Their dense smoothed synthesis and the
actual ambient factorization justify the traces before
their finite-dimensional values are summed.
The available cote~1 locators remain scans~079--080 and
111--112, ALPOGE\_BRIDGES BR-TR-01/02; no new archive
reading or Artin-motive realization is claimed here.
